%% Generated by Sphinx.
\def\sphinxdocclass{report}
\documentclass[letterpaper,10pt,english]{sphinxmanual}
\ifdefined\pdfpxdimen
   \let\sphinxpxdimen\pdfpxdimen\else\newdimen\sphinxpxdimen
\fi \sphinxpxdimen=.75bp\relax
\ifdefined\pdfimageresolution
    \pdfimageresolution= \numexpr \dimexpr1in\relax/\sphinxpxdimen\relax
\fi
\newdimen\sphinxremdimen\sphinxremdimen = 10pt
%% let collapsible pdf bookmarks panel have high depth per default
\PassOptionsToPackage{bookmarksdepth=5}{hyperref}

\PassOptionsToPackage{booktabs}{sphinx}
\PassOptionsToPackage{colorrows}{sphinx}

\PassOptionsToPackage{warn}{textcomp}
\usepackage[utf8]{inputenc}
\ifdefined\DeclareUnicodeCharacter
% support both utf8 and utf8x syntaxes
  \ifdefined\DeclareUnicodeCharacterAsOptional
    \def\sphinxDUC#1{\DeclareUnicodeCharacter{"#1}}
  \else
    \let\sphinxDUC\DeclareUnicodeCharacter
  \fi
  \sphinxDUC{00A0}{\nobreakspace}
  \sphinxDUC{2500}{\sphinxunichar{2500}}
  \sphinxDUC{2502}{\sphinxunichar{2502}}
  \sphinxDUC{2514}{\sphinxunichar{2514}}
  \sphinxDUC{251C}{\sphinxunichar{251C}}
  \sphinxDUC{2572}{\textbackslash}
\fi
\usepackage{cmap}
\usepackage[T1]{fontenc}
\usepackage{amsmath,amssymb,amstext}
\usepackage{babel}



\usepackage{tgtermes}
\usepackage{tgheros}
\renewcommand{\ttdefault}{txtt}



\usepackage[Bjarne]{fncychap}
\usepackage{sphinx}

\fvset{fontsize=auto}
\usepackage{geometry}


% Include hyperref last.
\usepackage{hyperref}
% Fix anchor placement for figures with captions.
\usepackage{hypcap}% it must be loaded after hyperref.
% Set up styles of URL: it should be placed after hyperref.
\urlstyle{same}


\usepackage{sphinxmessages}




\title{SAFE\sphinxhyphen{}M\sphinxhyphen{}ORP}
\date{Jul 03, 2026}
\release{1}
\author{Mathis, François Méthivier}
\newcommand{\sphinxlogo}{\vbox{}}
\renewcommand{\releasename}{Release}
\makeindex
\begin{document}

\ifdefined\shorthandoff
  \ifnum\catcode`\=\string=\active\shorthandoff{=}\fi
  \ifnum\catcode`\"=\active\shorthandoff{"}\fi
\fi

\pagestyle{empty}
\sphinxmaketitle
\pagestyle{plain}
\sphinxtableofcontents
\pagestyle{normal}
\phantomsection\label{\detokenize{index::doc}}


\sphinxAtStartPar
Add your content using \sphinxcode{\sphinxupquote{reStructuredText}} syntax. See the
\sphinxhref{https://www.sphinx-doc.org/en/master/usage/restructuredtext/index.html}{reStructuredText}
documentation for details.

\sphinxstepscope


\chapter{SRC}
\label{\detokenize{modules:src}}\label{\detokenize{modules::doc}}
\sphinxstepscope


\section{Mesure module}
\label{\detokenize{Mesure:module-Mesure}}\label{\detokenize{Mesure:mesure-module}}\label{\detokenize{Mesure::doc}}\index{module@\spxentry{module}!Mesure@\spxentry{Mesure}}\index{Mesure@\spxentry{Mesure}!module@\spxentry{module}}\index{Graphe\_T\_V() (in module Mesure)@\spxentry{Graphe\_T\_V()}\spxextra{in module Mesure}}

\begin{fulllineitems}
\phantomsection\label{\detokenize{Mesure:Mesure.Graphe_T_V}}
\pysigstartsignatures
\pysiglinewithargsret
{\sphinxcode{\sphinxupquote{Mesure.}}\sphinxbfcode{\sphinxupquote{Graphe\_T\_V}}}
{\sphinxparam{\DUrole{n}{x}}\sphinxparamcomma \sphinxparam{\DUrole{n}{y}}\sphinxparamcomma \sphinxparam{\DUrole{n}{t\_temps}}}
{}
\pysigstopsignatures
\sphinxAtStartPar
\_summary\_
Affiche deux graphiques en nuage de points des tensions mesurées (corrigées en mV) et de températures(°C)
Crée une nouvelle fenêtre matplotlib intitulée ‘Graphique Voltage et Température et trace
chaque valeur de y en fonction de son indice (axe temporel implicite).
\begin{quote}\begin{description}
\sphinxlineitem{Parameters}\begin{itemize}
\item {} 
\sphinxAtStartPar
\sphinxstyleliteralstrong{\sphinxupquote{x}} (\sphinxstyleliteralemphasis{\sphinxupquote{list}}\sphinxstyleliteralemphasis{\sphinxupquote{{[}}}\sphinxstyleliteralemphasis{\sphinxupquote{float}}\sphinxstyleliteralemphasis{\sphinxupquote{{]}}}) \textendash{} Liste des températures mesurées à afficher (en °C).

\item {} 
\sphinxAtStartPar
\sphinxstyleliteralstrong{\sphinxupquote{y}} (\sphinxstyleliteralemphasis{\sphinxupquote{list}}\sphinxstyleliteralemphasis{\sphinxupquote{{[}}}\sphinxstyleliteralemphasis{\sphinxupquote{float}}\sphinxstyleliteralemphasis{\sphinxupquote{{]}}}) \textendash{} Liste des tensions corrigées à afficher (en V ou mV selon calibration).

\item {} 
\sphinxAtStartPar
\sphinxstyleliteralstrong{\sphinxupquote{t\_temps}} (\sphinxstyleliteralemphasis{\sphinxupquote{list}}\sphinxstyleliteralemphasis{\sphinxupquote{{[}}}\sphinxstyleliteralemphasis{\sphinxupquote{float}}\sphinxstyleliteralemphasis{\sphinxupquote{{]}}}) \textendash{} Liste des temps écoulés depuis le début (en sec).

\end{itemize}

\end{description}\end{quote}

\end{fulllineitems}

\index{connexion\_port() (in module Mesure)@\spxentry{connexion\_port()}\spxextra{in module Mesure}}

\begin{fulllineitems}
\phantomsection\label{\detokenize{Mesure:Mesure.connexion_port}}
\pysigstartsignatures
\pysiglinewithargsret
{\sphinxcode{\sphinxupquote{Mesure.}}\sphinxbfcode{\sphinxupquote{connexion\_port}}}
{\sphinxparam{\DUrole{n}{br}\DUrole{o}{=}\DUrole{default_value}{115200}}\sphinxparamcomma \sphinxparam{\DUrole{n}{portIN}\DUrole{o}{=}\DUrole{default_value}{\textquotesingle{}\textquotesingle{}}}}
{}
\pysigstopsignatures
\sphinxAtStartPar
\_summary\_
Établit la connexion série avec l’Arduino. Si un port est fourni manuellement,
tente de s’y connecter directement. Sinon, détecte automatiquement le port
Arduino parmi les ports disponibles (Windows et Linux), en cherchant les
identifiants caractéristiques dans le nom du port.
\begin{quote}\begin{description}
\sphinxlineitem{Parameters}\begin{itemize}
\item {} 
\sphinxAtStartPar
\sphinxstyleliteralstrong{\sphinxupquote{br}} (\sphinxstyleliteralemphasis{\sphinxupquote{int}}) \textendash{} Baudrate de la connexion série. Par défaut 115200.

\item {} 
\sphinxAtStartPar
\sphinxstyleliteralstrong{\sphinxupquote{portIN}} (\sphinxstyleliteralemphasis{\sphinxupquote{str}}) \textendash{} Nom du port série à utiliser. Si vide, détection automatique.

\end{itemize}

\sphinxlineitem{Returns}
\sphinxAtStartPar
\begin{itemize}
\item {} 
\sphinxAtStartPar
portIN : nom du port détecté ou fourni (ex. ‘COM6’, ‘/dev/ttyACM0’).

\item {} 
\sphinxAtStartPar
s : objet serial.Serial connecté, ou la chaîne ‘error’ en cas d’échec.

\end{itemize}


\sphinxlineitem{Return type}
\sphinxAtStartPar
tuple{[}str, serial.Serial | str{]}

\end{description}\end{quote}

\end{fulllineitems}

\index{data() (in module Mesure)@\spxentry{data()}\spxextra{in module Mesure}}

\begin{fulllineitems}
\phantomsection\label{\detokenize{Mesure:Mesure.data}}
\pysigstartsignatures
\pysiglinewithargsret
{\sphinxcode{\sphinxupquote{Mesure.}}\sphinxbfcode{\sphinxupquote{data}}}
{\sphinxparam{\DUrole{n}{s}}\sphinxparamcomma \sphinxparam{\DUrole{n}{N}\DUrole{o}{=}\DUrole{default_value}{None}}}
{}
\pysigstopsignatures
\sphinxAtStartPar
\_summary\_
Acquiert un nombre fixe de mesures depuis le port série, espacées d’un
intervalle de temps connu (time\_inter), et renvoie un temps écoulé en
secondes pour chaque mesure (plutôt qu’un simple indice d’acquisition).
\begin{quote}\begin{description}
\sphinxlineitem{Parameters}\begin{itemize}
\item {} 
\sphinxAtStartPar
\sphinxstyleliteralstrong{\sphinxupquote{s}} (\sphinxstyleliteralemphasis{\sphinxupquote{serial.Serial}}) \textendash{} Objet Serial ouvert sur lequel lire les mesures.

\item {} 
\sphinxAtStartPar
\sphinxstyleliteralstrong{\sphinxupquote{N}} (\sphinxstyleliteralemphasis{\sphinxupquote{int}}\sphinxstyleliteralemphasis{\sphinxupquote{, }}\sphinxstyleliteralemphasis{\sphinxupquote{optional}}) \textendash{} Nombre de mesures à effectuer. Si None, demandé à l’utilisateur.

\item {} 
\sphinxAtStartPar
\sphinxstyleliteralstrong{\sphinxupquote{time\_inter}} (\sphinxstyleliteralemphasis{\sphinxupquote{float}}\sphinxstyleliteralemphasis{\sphinxupquote{, }}\sphinxstyleliteralemphasis{\sphinxupquote{optional}}) \textendash{} Intervalle de temps (en sec) entre deux lectures. Par défaut 0.08 sec.

\item {} 
\sphinxAtStartPar
\sphinxstyleliteralstrong{\sphinxupquote{Returns}} \textendash{} \begin{description}
\sphinxlineitem{tuple{[}list{[}float{]}, list{[}float{]}, list{[}float{]}{]}:}\begin{itemize}
\item {} 
\sphinxAtStartPar
T : liste des températures mesurées (en °C).

\item {} 
\sphinxAtStartPar
V : liste des tensions mesurées (en V).

\item {} 
\sphinxAtStartPar
t\_temps : liste des temps écoulés depuis le début (en sec).

\end{itemize}

\end{description}


\end{itemize}

\end{description}\end{quote}

\end{fulllineitems}

\index{enregistrement\_csv() (in module Mesure)@\spxentry{enregistrement\_csv()}\spxextra{in module Mesure}}

\begin{fulllineitems}
\phantomsection\label{\detokenize{Mesure:Mesure.enregistrement_csv}}
\pysigstartsignatures
\pysiglinewithargsret
{\sphinxcode{\sphinxupquote{Mesure.}}\sphinxbfcode{\sphinxupquote{enregistrement\_csv}}}
{\sphinxparam{\DUrole{n}{T}}\sphinxparamcomma \sphinxparam{\DUrole{n}{V\_real}}\sphinxparamcomma \sphinxparam{\DUrole{n}{moy}\DUrole{o}{=}\DUrole{default_value}{None}}\sphinxparamcomma \sphinxparam{\DUrole{n}{sigma}\DUrole{o}{=}\DUrole{default_value}{None}}}
{}
\pysigstopsignatures
\sphinxAtStartPar
\_summary\_
Enregistre les données de température et de potentiel calibré dans un fichier
CSV horodaté, accompagnées de la moyenne et de l’écart\sphinxhyphen{}type des tensions.
Le fichier est sauvegardé dans le dossier Data/data\_(T,V)/.
\begin{quote}\begin{description}
\sphinxlineitem{Parameters}\begin{itemize}
\item {} 
\sphinxAtStartPar
\sphinxstyleliteralstrong{\sphinxupquote{T}} (\sphinxstyleliteralemphasis{\sphinxupquote{list}}\sphinxstyleliteralemphasis{\sphinxupquote{{[}}}\sphinxstyleliteralemphasis{\sphinxupquote{float}}\sphinxstyleliteralemphasis{\sphinxupquote{{]}}}) \textendash{} Liste des températures mesurées (en °C).

\item {} 
\sphinxAtStartPar
\sphinxstyleliteralstrong{\sphinxupquote{V\_real}} (\sphinxstyleliteralemphasis{\sphinxupquote{list}}\sphinxstyleliteralemphasis{\sphinxupquote{{[}}}\sphinxstyleliteralemphasis{\sphinxupquote{float}}\sphinxstyleliteralemphasis{\sphinxupquote{{]}}}) \textendash{} Liste des potentiels calibrés (en mV).

\item {} 
\sphinxAtStartPar
\sphinxstyleliteralstrong{\sphinxupquote{moy}} (\sphinxstyleliteralemphasis{\sphinxupquote{float}}\sphinxstyleliteralemphasis{\sphinxupquote{, }}\sphinxstyleliteralemphasis{\sphinxupquote{optional}}) \textendash{} Moyenne des tensions calibrées. Calculée automatiquement
si non fournie.

\item {} 
\sphinxAtStartPar
\sphinxstyleliteralstrong{\sphinxupquote{sigma}} (\sphinxstyleliteralemphasis{\sphinxupquote{float}}\sphinxstyleliteralemphasis{\sphinxupquote{, }}\sphinxstyleliteralemphasis{\sphinxupquote{optional}}) \textendash{} Écart\sphinxhyphen{}type des tensions calibrées. Calculé
automatiquement si non fourni.

\end{itemize}

\sphinxlineitem{Returns}
\sphinxAtStartPar
Message de confirmation d’enregistrement.

\sphinxlineitem{Return type}
\sphinxAtStartPar
str

\end{description}\end{quote}

\end{fulllineitems}

\index{enregistrement\_pdf() (in module Mesure)@\spxentry{enregistrement\_pdf()}\spxextra{in module Mesure}}

\begin{fulllineitems}
\phantomsection\label{\detokenize{Mesure:Mesure.enregistrement_pdf}}
\pysigstartsignatures
\pysiglinewithargsret
{\sphinxcode{\sphinxupquote{Mesure.}}\sphinxbfcode{\sphinxupquote{enregistrement\_pdf}}}
{\sphinxparam{\DUrole{n}{figure}}}
{}
\pysigstopsignatures
\sphinxAtStartPar
\_summary\_
Enregistre le graphique de température et de potentiel sous forme d’image PDF
horodatée dans le dossier Data/data\_figures/data\_figures\_mesures/.
\begin{quote}\begin{description}
\sphinxlineitem{Parameters}
\sphinxAtStartPar
\sphinxstyleliteralstrong{\sphinxupquote{figure}} (\sphinxstyleliteralemphasis{\sphinxupquote{matplotlib.figure.Figure}}) \textendash{} Figure matplotlib à sauvegarder.

\sphinxlineitem{Returns}
\sphinxAtStartPar
Message de confirmation d’enregistrement.

\sphinxlineitem{Return type}
\sphinxAtStartPar
str

\end{description}\end{quote}

\end{fulllineitems}

\index{graphe\_csv() (in module Mesure)@\spxentry{graphe\_csv()}\spxextra{in module Mesure}}

\begin{fulllineitems}
\phantomsection\label{\detokenize{Mesure:Mesure.graphe_csv}}
\pysigstartsignatures
\pysiglinewithargsret
{\sphinxcode{\sphinxupquote{Mesure.}}\sphinxbfcode{\sphinxupquote{graphe\_csv}}}
{\sphinxparam{\DUrole{n}{chemin}}}
{}
\pysigstopsignatures
\sphinxAtStartPar
\_summary\_
Lit un fichier CSV contenant des données de température et de potentiel, puis
trace un graphique de T et V en fonction du nombre de mesures.
Le fichier CSV doit être horodaté et contenir deux colonnes : température (°C) et potentiel (mV).

\end{fulllineitems}

\index{graphe\_live() (in module Mesure)@\spxentry{graphe\_live()}\spxextra{in module Mesure}}

\begin{fulllineitems}
\phantomsection\label{\detokenize{Mesure:Mesure.graphe_live}}
\pysigstartsignatures
\pysiglinewithargsret
{\sphinxcode{\sphinxupquote{Mesure.}}\sphinxbfcode{\sphinxupquote{graphe\_live}}}
{\sphinxparam{\DUrole{n}{C0}}\sphinxparamcomma \sphinxparam{\DUrole{n}{s}}}
{}
\pysigstopsignatures
\end{fulllineitems}

\index{informations\_1er\_ordre() (in module Mesure)@\spxentry{informations\_1er\_ordre()}\spxextra{in module Mesure}}

\begin{fulllineitems}
\phantomsection\label{\detokenize{Mesure:Mesure.informations_1er_ordre}}
\pysigstartsignatures
\pysiglinewithargsret
{\sphinxcode{\sphinxupquote{Mesure.}}\sphinxbfcode{\sphinxupquote{informations\_1er\_ordre}}}
{\sphinxparam{\DUrole{n}{T}}\sphinxparamcomma \sphinxparam{\DUrole{n}{V}}\sphinxparamcomma \sphinxparam{\DUrole{n}{a}}}
{}
\pysigstopsignatures
\sphinxAtStartPar
\_summary\_
Calcule et affiche les statistiques de base d’une session de mesures : moyenne et écart\sphinxhyphen{}type
des tensions calibrées et des températures mesurées. Affiche les résultats dans la console.
\begin{quote}\begin{description}
\sphinxlineitem{Parameters}\begin{itemize}
\item {} 
\sphinxAtStartPar
\sphinxstyleliteralstrong{\sphinxupquote{T}} (\sphinxstyleliteralemphasis{\sphinxupquote{list}}\sphinxstyleliteralemphasis{\sphinxupquote{{[}}}\sphinxstyleliteralemphasis{\sphinxupquote{float}}\sphinxstyleliteralemphasis{\sphinxupquote{{]}}}) \textendash{} Liste des températures mesurées (en °C).

\item {} 
\sphinxAtStartPar
\sphinxstyleliteralstrong{\sphinxupquote{V}} (\sphinxstyleliteralemphasis{\sphinxupquote{list}}\sphinxstyleliteralemphasis{\sphinxupquote{{[}}}\sphinxstyleliteralemphasis{\sphinxupquote{float}}\sphinxstyleliteralemphasis{\sphinxupquote{{]}}}) \textendash{} Liste des tensions calibrées (en mV).

\end{itemize}

\end{description}\end{quote}

\end{fulllineitems}

\index{t\_attente() (in module Mesure)@\spxentry{t\_attente()}\spxextra{in module Mesure}}

\begin{fulllineitems}
\phantomsection\label{\detokenize{Mesure:Mesure.t_attente}}
\pysigstartsignatures
\pysiglinewithargsret
{\sphinxcode{\sphinxupquote{Mesure.}}\sphinxbfcode{\sphinxupquote{t\_attente}}}
{\sphinxparam{\DUrole{n}{nb}\DUrole{o}{=}\DUrole{default_value}{None}}}
{}
\pysigstopsignatures
\end{fulllineitems}


\sphinxstepscope


\section{calibration module}
\label{\detokenize{calibration:module-calibration}}\label{\detokenize{calibration:calibration-module}}\label{\detokenize{calibration::doc}}\index{module@\spxentry{module}!calibration@\spxentry{calibration}}\index{calibration@\spxentry{calibration}!module@\spxentry{module}}\index{V\_real\_f() (in module calibration)@\spxentry{V\_real\_f()}\spxextra{in module calibration}}

\begin{fulllineitems}
\phantomsection\label{\detokenize{calibration:calibration.V_real_f}}
\pysigstartsignatures
\pysiglinewithargsret
{\sphinxcode{\sphinxupquote{calibration.}}\sphinxbfcode{\sphinxupquote{V\_real\_f}}}
{\sphinxparam{\DUrole{n}{V}}\sphinxparamcomma \sphinxparam{\DUrole{n}{C0}}}
{}
\pysigstopsignatures
\sphinxAtStartPar
\_summary\_
Convertit une liste de tensions brutes en valeurs de potentiel corrigées (en mV),
en appliquant la formule de conversion du capteur et l’offset de calibration C0.
La conversion utilisée est : V\_corrigé = (2 \sphinxhyphen{} i) * 1000 + C0.
\begin{quote}\begin{description}
\sphinxlineitem{Parameters}\begin{itemize}
\item {} 
\sphinxAtStartPar
\sphinxstyleliteralstrong{\sphinxupquote{V}} (\sphinxstyleliteralemphasis{\sphinxupquote{list}}\sphinxstyleliteralemphasis{\sphinxupquote{{[}}}\sphinxstyleliteralemphasis{\sphinxupquote{float}}\sphinxstyleliteralemphasis{\sphinxupquote{{]}}}) \textendash{} Liste des tensions brutes mesurées par le capteur (en volts).

\item {} 
\sphinxAtStartPar
\sphinxstyleliteralstrong{\sphinxupquote{C0}} (\sphinxstyleliteralemphasis{\sphinxupquote{float}}) \textendash{} Terme correctif issu de la calibration (en mV), peut être = 0 si
calibration d’usine, ou une valeur calculée via calibration\_etalon().

\end{itemize}

\sphinxlineitem{Returns}
\sphinxAtStartPar
Liste des potentiels corrigés (en mV), de même longueur que V.

\sphinxlineitem{Return type}
\sphinxAtStartPar
list{[}float{]}

\end{description}\end{quote}

\end{fulllineitems}

\index{calibration\_2\_etalons() (in module calibration)@\spxentry{calibration\_2\_etalons()}\spxextra{in module calibration}}

\begin{fulllineitems}
\phantomsection\label{\detokenize{calibration:calibration.calibration_2_etalons}}
\pysigstartsignatures
\pysiglinewithargsret
{\sphinxcode{\sphinxupquote{calibration.}}\sphinxbfcode{\sphinxupquote{calibration\_2\_etalons}}}
{\sphinxparam{\DUrole{n}{s}}\sphinxparamcomma \sphinxparam{\DUrole{n}{E1}\DUrole{o}{=}\DUrole{default_value}{None}}\sphinxparamcomma \sphinxparam{\DUrole{n}{E2}\DUrole{o}{=}\DUrole{default_value}{None}}\sphinxparamcomma \sphinxparam{\DUrole{n}{V1}\DUrole{o}{=}\DUrole{default_value}{None}}\sphinxparamcomma \sphinxparam{\DUrole{n}{V2}\DUrole{o}{=}\DUrole{default_value}{None}}}
{}
\pysigstopsignatures
\sphinxAtStartPar
\_summary\_
Effectue une calibration à deux étalons par régression linéaire.
Demande à l’utilisateur les potentiels des deux solutions étalons, acquiert
les mesures sur chacune d’elles, puis calcule la droite de tendance par
régression linéaire (scipy.stats.linregress) entre les tensions moyennes
mesurées et les potentiels théoriques. Avertit si le R\(\sp{\text{2}}\) est inférieur à 0.9.

\sphinxAtStartPar
Remarque : avec seulement 2 points, R\(\sp{\text{2}}\) vaut toujours 1 par construction
mathématique — l’avertissement R\(\sp{\text{2}}\) \textless{} 0.9 ne peut donc jamais se déclencher.
\begin{quote}\begin{description}
\sphinxlineitem{Parameters}\begin{itemize}
\item {} 
\sphinxAtStartPar
\sphinxstyleliteralstrong{\sphinxupquote{s}} (\sphinxstyleliteralemphasis{\sphinxupquote{serial.Serial}}) \textendash{} Objet de connexion série avec l’Arduino.

\item {} 
\sphinxAtStartPar
\sphinxstyleliteralstrong{\sphinxupquote{E1}} (\sphinxstyleliteralemphasis{\sphinxupquote{float}}\sphinxstyleliteralemphasis{\sphinxupquote{, }}\sphinxstyleliteralemphasis{\sphinxupquote{optional}}) \textendash{} Potentiel théorique de la solution étalon 1 (en mV).

\item {} 
\sphinxAtStartPar
\sphinxstyleliteralstrong{\sphinxupquote{E2}} (\sphinxstyleliteralemphasis{\sphinxupquote{float}}\sphinxstyleliteralemphasis{\sphinxupquote{, }}\sphinxstyleliteralemphasis{\sphinxupquote{optional}}) \textendash{} Potentiel théorique de la solution étalon 2 (en mV).

\item {} 
\sphinxAtStartPar
\sphinxstyleliteralstrong{\sphinxupquote{V1}} (\sphinxstyleliteralemphasis{\sphinxupquote{list}}\sphinxstyleliteralemphasis{\sphinxupquote{{[}}}\sphinxstyleliteralemphasis{\sphinxupquote{float}}\sphinxstyleliteralemphasis{\sphinxupquote{{]}}}\sphinxstyleliteralemphasis{\sphinxupquote{, }}\sphinxstyleliteralemphasis{\sphinxupquote{optional}}) \textendash{} Tensions brutes mesurées sur l’étalon 1.

\item {} 
\sphinxAtStartPar
\sphinxstyleliteralstrong{\sphinxupquote{V2}} (\sphinxstyleliteralemphasis{\sphinxupquote{list}}\sphinxstyleliteralemphasis{\sphinxupquote{{[}}}\sphinxstyleliteralemphasis{\sphinxupquote{float}}\sphinxstyleliteralemphasis{\sphinxupquote{{]}}}\sphinxstyleliteralemphasis{\sphinxupquote{, }}\sphinxstyleliteralemphasis{\sphinxupquote{optional}}) \textendash{} Tensions brutes mesurées sur l’étalon 2.

\item {} 
\sphinxAtStartPar
\sphinxstyleliteralstrong{\sphinxupquote{V\_real1}} \textendash{} tension convertit en mV

\item {} 
\sphinxAtStartPar
\sphinxstyleliteralstrong{\sphinxupquote{V\_real2}} \textendash{} tension convertit en mV

\end{itemize}

\end{description}\end{quote}

\sphinxAtStartPar
E1 et E2 sont demandé à l’utilisateur si non fournis, V1 et V2 sont acquises via mes.data() si non fournies
V\_real1 et V\_real2 sont acquise via V1 et V2
\begin{quote}\begin{description}
\sphinxlineitem{Returns}
\sphinxAtStartPar
\begin{description}
\sphinxlineitem{(C0, V1\_moy, V2\_moy, tendance, r\_squared, E1, E2)}\begin{itemize}
\item {} 
\sphinxAtStartPar
C0 (float): Intercept de la droite, utilisé comme offset de calibration (en mV).

\item {} 
\sphinxAtStartPar
V1\_moy (float): Tension moyenne mesurée sur l’étalon 1 (en mV).

\item {} 
\sphinxAtStartPar
V2\_moy (float): Tension moyenne mesurée sur l’étalon 2 (en mV).

\item {} 
\sphinxAtStartPar
tendance (np.poly1d): Droite de calibration sous forme de polynôme.

\item {} 
\sphinxAtStartPar
r\_squared (float): Coefficient de détermination R\(\sp{\text{2}}\) de la régression.

\item {} 
\sphinxAtStartPar
E1 (float): Potentiel théorique de l’étalon 1 (en mV).

\item {} 
\sphinxAtStartPar
E2 (float): Potentiel théorique de l’étalon 2 (en mV).

\end{itemize}

\end{description}


\sphinxlineitem{Return type}
\sphinxAtStartPar
tuple

\end{description}\end{quote}

\end{fulllineitems}

\index{calibration\_etalon() (in module calibration)@\spxentry{calibration\_etalon()}\spxextra{in module calibration}}

\begin{fulllineitems}
\phantomsection\label{\detokenize{calibration:calibration.calibration_etalon}}
\pysigstartsignatures
\pysiglinewithargsret
{\sphinxcode{\sphinxupquote{calibration.}}\sphinxbfcode{\sphinxupquote{calibration\_etalon}}}
{\sphinxparam{\DUrole{n}{V}}}
{}
\pysigstopsignatures
\sphinxAtStartPar
\_summary\_
Calcule le terme correctif C0 par calibration à un étalon.
Demande à l’utilisateur la valeur du potentiel de la solution étalon,
puis calcule pour chaque tension mesurée l’écart entre la valeur attendue
et la valeur convertie. C0 est la moyenne de ces écarts, utilisée ensuite
pour corriger les mesures réelles.
\begin{quote}\begin{description}
\sphinxlineitem{Parameters}
\sphinxAtStartPar
\sphinxstyleliteralstrong{\sphinxupquote{V}} (\sphinxstyleliteralemphasis{\sphinxupquote{list}}\sphinxstyleliteralemphasis{\sphinxupquote{{[}}}\sphinxstyleliteralemphasis{\sphinxupquote{float}}\sphinxstyleliteralemphasis{\sphinxupquote{{]}}}) \textendash{} Liste des tensions brutes mesurées sur la solution étalon (en volts).

\sphinxlineitem{Returns}
\sphinxAtStartPar
C0, terme correctif moyen (offset en mV) à appliquer aux mesures futures.

\sphinxlineitem{Return type}
\sphinxAtStartPar
float

\end{description}\end{quote}

\end{fulllineitems}

\index{calibration\_usine() (in module calibration)@\spxentry{calibration\_usine()}\spxextra{in module calibration}}

\begin{fulllineitems}
\phantomsection\label{\detokenize{calibration:calibration.calibration_usine}}
\pysigstartsignatures
\pysiglinewithargsret
{\sphinxcode{\sphinxupquote{calibration.}}\sphinxbfcode{\sphinxupquote{calibration\_usine}}}
{}
{}
\pysigstopsignatures\begin{description}
\sphinxlineitem{\_summary\_}
\sphinxAtStartPar
Applique la formule de calibration d’usine à une liste de tensions mesurées.
Pour chaque valeur de V (‘brut’), on convertit : E = (2 \sphinxhyphen{} i) * 1000 (qui sera calculé par V\_real), mais retourne
systématiquement C0 = 0 (aucun offset correctif n’est appliqué).
Cette fonction sert de référence sans correction : elle suppose que le capteur
est déjà calibré en sortie d’usine.

\end{description}

\sphinxAtStartPar
Remarque : on ne savait pas si on devait définir une fonction (cacalibration\_usine) ou  juste mettre C0 = 0 pour l’option 1 du menu
\begin{quote}
\begin{description}
\sphinxlineitem{Args:}
\sphinxAtStartPar
none

\sphinxlineitem{Returns:}
\sphinxAtStartPar
int: C0 = 0, terme correctif nul (calibration d’usine sans offset).

\end{description}
\end{quote}

\end{fulllineitems}

\index{enregistrement\_cal() (in module calibration)@\spxentry{enregistrement\_cal()}\spxextra{in module calibration}}

\begin{fulllineitems}
\phantomsection\label{\detokenize{calibration:calibration.enregistrement_cal}}
\pysigstartsignatures
\pysiglinewithargsret
{\sphinxcode{\sphinxupquote{calibration.}}\sphinxbfcode{\sphinxupquote{enregistrement\_cal}}}
{\sphinxparam{\DUrole{n}{C0}}\sphinxparamcomma \sphinxparam{\DUrole{n}{tendance}\DUrole{o}{=}\DUrole{default_value}{None}}\sphinxparamcomma \sphinxparam{\DUrole{n}{nb\_etalons}\DUrole{o}{=}\DUrole{default_value}{None}}}
{}
\pysigstopsignatures
\sphinxAtStartPar
\_summary\_
Enregistre les paramètres de calibration dans un fichier CSV horodaté.
Selon le nombre d’étalons, le fichier contient soit uniquement C0 (calibration
à 1 étalon ou d’usine), soit C0, la pente et l’intercept de la droite de
tendance (calibration à 2 étalons). Le fichier est sauvegardé dans le dossier
Data/data\_calibration/.
\begin{quote}\begin{description}
\sphinxlineitem{Parameters}\begin{itemize}
\item {} 
\sphinxAtStartPar
\sphinxstyleliteralstrong{\sphinxupquote{C0}} (\sphinxstyleliteralemphasis{\sphinxupquote{float}}) \textendash{} Terme correctif issu de la calibration (en mV).

\item {} 
\sphinxAtStartPar
\sphinxstyleliteralstrong{\sphinxupquote{tendance}} (\sphinxstyleliteralemphasis{\sphinxupquote{np.poly1d}}\sphinxstyleliteralemphasis{\sphinxupquote{, }}\sphinxstyleliteralemphasis{\sphinxupquote{optional}}) \textendash{} Objet polynôme contenant la pente et
l’intercept de la droite d’étalonnage. None si calibration à 1 étalon.

\item {} 
\sphinxAtStartPar
\sphinxstyleliteralstrong{\sphinxupquote{nb\_etalons}} (\sphinxstyleliteralemphasis{\sphinxupquote{int}}\sphinxstyleliteralemphasis{\sphinxupquote{, }}\sphinxstyleliteralemphasis{\sphinxupquote{optional}}) \textendash{} Nombre d’étalons utilisés (1 ou 2). None si
calibration d’usine.

\end{itemize}

\sphinxlineitem{Returns}
\sphinxAtStartPar
None

\end{description}\end{quote}

\end{fulllineitems}

\index{enregistrement\_cal2\_pdf() (in module calibration)@\spxentry{enregistrement\_cal2\_pdf()}\spxextra{in module calibration}}

\begin{fulllineitems}
\phantomsection\label{\detokenize{calibration:calibration.enregistrement_cal2_pdf}}
\pysigstartsignatures
\pysiglinewithargsret
{\sphinxcode{\sphinxupquote{calibration.}}\sphinxbfcode{\sphinxupquote{enregistrement\_cal2\_pdf}}}
{\sphinxparam{\DUrole{n}{fig}}}
{}
\pysigstopsignatures
\sphinxAtStartPar
\_summary\_
Enregistre le graphique de calibration à 2 étalons sous forme d’image PDF
horodatée dans le dossier Data/data\_figures/data\_figures\_calibration/.
\begin{quote}\begin{description}
\sphinxlineitem{Parameters}
\sphinxAtStartPar
\sphinxstyleliteralstrong{\sphinxupquote{fig}} (\sphinxstyleliteralemphasis{\sphinxupquote{matplotlib.figure.Figure}}) \textendash{} Figure matplotlib à sauvegarder.

\sphinxlineitem{Returns}
\sphinxAtStartPar
Message de confirmation d’enregistrement.

\sphinxlineitem{Return type}
\sphinxAtStartPar
str

\end{description}\end{quote}

\end{fulllineitems}

\index{graphe\_cal2() (in module calibration)@\spxentry{graphe\_cal2()}\spxextra{in module calibration}}

\begin{fulllineitems}
\phantomsection\label{\detokenize{calibration:calibration.graphe_cal2}}
\pysigstartsignatures
\pysiglinewithargsret
{\sphinxcode{\sphinxupquote{calibration.}}\sphinxbfcode{\sphinxupquote{graphe\_cal2}}}
{\sphinxparam{\DUrole{n}{tendance}}\sphinxparamcomma \sphinxparam{\DUrole{n}{a}}\sphinxparamcomma \sphinxparam{\DUrole{n}{C0}}\sphinxparamcomma \sphinxparam{\DUrole{n}{E1}}\sphinxparamcomma \sphinxparam{\DUrole{n}{E2}}\sphinxparamcomma \sphinxparam{\DUrole{n}{V1\_real}}\sphinxparamcomma \sphinxparam{\DUrole{n}{V2\_real}}}
{}
\pysigstopsignatures
\sphinxAtStartPar
\_\_summary\_\_
Affiche le graphique de la droite de calibration à 2 étalons.
Trace la droite de calibration E = a * V\_real + C0 ainsi que les deux points
étalons, avec l’équation de la droite en légende.
\begin{quote}\begin{description}
\sphinxlineitem{Parameters}\begin{itemize}
\item {} 
\sphinxAtStartPar
\sphinxstyleliteralstrong{\sphinxupquote{tendance}} (\sphinxstyleliteralemphasis{\sphinxupquote{np.poly1d}}) \textendash{} Droite de calibration sous forme de polynôme,
construite à partir de a et C0 via np.poly1d({[}a, C0{]}).

\item {} 
\sphinxAtStartPar
\sphinxstyleliteralstrong{\sphinxupquote{a}} (\sphinxstyleliteralemphasis{\sphinxupquote{float}}) \textendash{} Pente de la droite de calibration (sans unité, idéalement ≈ 1).

\item {} 
\sphinxAtStartPar
\sphinxstyleliteralstrong{\sphinxupquote{C0}} (\sphinxstyleliteralemphasis{\sphinxupquote{float}}) \textendash{} Intercept de la droite de calibration (en mV).

\item {} 
\sphinxAtStartPar
\sphinxstyleliteralstrong{\sphinxupquote{E1}} (\sphinxstyleliteralemphasis{\sphinxupquote{float}}) \textendash{} Potentiel théorique de la solution étalon 1 (en mV).

\item {} 
\sphinxAtStartPar
\sphinxstyleliteralstrong{\sphinxupquote{E2}} (\sphinxstyleliteralemphasis{\sphinxupquote{float}}) \textendash{} Potentiel théorique de la solution étalon 2 (en mV).

\item {} 
\sphinxAtStartPar
\sphinxstyleliteralstrong{\sphinxupquote{V1\_real}} (\sphinxstyleliteralemphasis{\sphinxupquote{float}}) \textendash{} Tension mesurée convertie de l’étalon 1 (en mV),
obtenue via (2 \sphinxhyphen{} V1\_moy) * 1000.

\item {} 
\sphinxAtStartPar
\sphinxstyleliteralstrong{\sphinxupquote{V2\_real}} (\sphinxstyleliteralemphasis{\sphinxupquote{float}}) \textendash{} Tension mesurée convertie de l’étalon 2 (en mV),
obtenue via (2 \sphinxhyphen{} V2\_moy) * 1000.

\end{itemize}

\sphinxlineitem{Returns}
\sphinxAtStartPar
Figure matplotlib du graphique de calibration.

\sphinxlineitem{Return type}
\sphinxAtStartPar
matplotlib.figure.Figure

\end{description}\end{quote}

\end{fulllineitems}


\sphinxstepscope


\section{consigne module}
\label{\detokenize{consigne:consigne-module}}\label{\detokenize{consigne::doc}}
\sphinxstepscope


\section{test module}
\label{\detokenize{test:module-test}}\label{\detokenize{test:test-module}}\label{\detokenize{test::doc}}\index{module@\spxentry{module}!test@\spxentry{test}}\index{test@\spxentry{test}!module@\spxentry{module}}\index{graphe\_cal2() (in module test)@\spxentry{graphe\_cal2()}\spxextra{in module test}}

\begin{fulllineitems}
\phantomsection\label{\detokenize{test:test.graphe_cal2}}
\pysigstartsignatures
\pysiglinewithargsret
{\sphinxcode{\sphinxupquote{test.}}\sphinxbfcode{\sphinxupquote{graphe\_cal2}}}
{\sphinxparam{\DUrole{n}{tendance}}\sphinxparamcomma \sphinxparam{\DUrole{n}{a}}\sphinxparamcomma \sphinxparam{\DUrole{n}{C0}}\sphinxparamcomma \sphinxparam{\DUrole{n}{E1}}\sphinxparamcomma \sphinxparam{\DUrole{n}{E2}}\sphinxparamcomma \sphinxparam{\DUrole{n}{V1\_real}}\sphinxparamcomma \sphinxparam{\DUrole{n}{V2\_real}}}
{}
\pysigstopsignatures
\sphinxAtStartPar
\_\_summary\_\_
Affiche le graphique de la droite de calibration à 2 étalons.
Trace la droite de calibration E = a * V\_real + C0 ainsi que les deux points
étalons, avec l’équation de la droite en légende.
\begin{quote}\begin{description}
\sphinxlineitem{Parameters}\begin{itemize}
\item {} 
\sphinxAtStartPar
\sphinxstyleliteralstrong{\sphinxupquote{tendance}} (\sphinxstyleliteralemphasis{\sphinxupquote{np.poly1d}}) \textendash{} Droite de calibration sous forme de polynôme,
construite à partir de a et C0 via np.poly1d({[}a, C0{]}).

\item {} 
\sphinxAtStartPar
\sphinxstyleliteralstrong{\sphinxupquote{a}} (\sphinxstyleliteralemphasis{\sphinxupquote{float}}) \textendash{} Pente de la droite de calibration (sans unité, idéalement ≈ 1).

\item {} 
\sphinxAtStartPar
\sphinxstyleliteralstrong{\sphinxupquote{C0}} (\sphinxstyleliteralemphasis{\sphinxupquote{float}}) \textendash{} Intercept de la droite de calibration (en mV).

\item {} 
\sphinxAtStartPar
\sphinxstyleliteralstrong{\sphinxupquote{E1}} (\sphinxstyleliteralemphasis{\sphinxupquote{float}}) \textendash{} Potentiel théorique de la solution étalon 1 (en mV).

\item {} 
\sphinxAtStartPar
\sphinxstyleliteralstrong{\sphinxupquote{E2}} (\sphinxstyleliteralemphasis{\sphinxupquote{float}}) \textendash{} Potentiel théorique de la solution étalon 2 (en mV).

\item {} 
\sphinxAtStartPar
\sphinxstyleliteralstrong{\sphinxupquote{V1\_real}} (\sphinxstyleliteralemphasis{\sphinxupquote{float}}) \textendash{} Tension mesurée convertie de l’étalon 1 (en mV),
obtenue via (2 \sphinxhyphen{} V1\_moy) * 1000.

\item {} 
\sphinxAtStartPar
\sphinxstyleliteralstrong{\sphinxupquote{V2\_real}} (\sphinxstyleliteralemphasis{\sphinxupquote{float}}) \textendash{} Tension mesurée convertie de l’étalon 2 (en mV),
obtenue via (2 \sphinxhyphen{} V2\_moy) * 1000.

\end{itemize}

\sphinxlineitem{Returns}
\sphinxAtStartPar
Figure matplotlib du graphique de calibration.

\sphinxlineitem{Return type}
\sphinxAtStartPar
matplotlib.figure.Figure

\end{description}\end{quote}

\end{fulllineitems}



\renewcommand{\indexname}{Python Module Index}
\begin{sphinxtheindex}
\let\bigletter\sphinxstyleindexlettergroup
\bigletter{c}
\item\relax\sphinxstyleindexentry{calibration}\sphinxstyleindexpageref{calibration:\detokenize{module-calibration}}
\indexspace
\bigletter{m}
\item\relax\sphinxstyleindexentry{Mesure}\sphinxstyleindexpageref{Mesure:\detokenize{module-Mesure}}
\indexspace
\bigletter{t}
\item\relax\sphinxstyleindexentry{test}\sphinxstyleindexpageref{test:\detokenize{module-test}}
\end{sphinxtheindex}

\renewcommand{\indexname}{Index}
\printindex
\end{document}